% ============================================================
%  Chương 1 & 2 — Đồ án tốt nghiệp
%  "Nghiên cứu, phát triển nền tảng mã hoá video cho các ứng dụng thị giác máy"
% ============================================================

\chapter{Đặt vấn đề}

\section{Bối cảnh và lý do chọn đề tài}

Trong khoảng hai thập kỷ trở lại đây, ngành công nghiệp ô tô đã chứng kiến một quá trình chuyển dịch sâu sắc từ phương tiện cơ giới truyền thống sang phương tiện thông minh có khả năng hỗ trợ và tự động hoá việc lái xe. Các hệ thống hỗ trợ lái nâng cao (Advanced Driver Assistance Systems) và lái xe tự động (Automated Driving) ngày càng được tích hợp rộng rãi trên các dòng xe thương mại, từ những tính năng đơn giản như cảnh báo lệch làn đường hay phanh khẩn cấp tự động cho đến những hệ thống phức tạp hơn có khả năng giữ làn, đỗ xe tự động và điều khiển hành trình thích ứng. Tổ chức Kỹ sư Ô tô Quốc tế (Society of Automotive Engineers) đã phân loại mức độ tự động hoá thành sáu cấp từ L0 đến L5, trong đó các phương tiện hiện đại trên thị trường chủ yếu đang dao động trong khoảng L1 đến L3, còn các cấp L4 và L5 vẫn đang là mục tiêu nghiên cứu của các hãng xe và phòng thí nghiệm hàng đầu thế giới.

Bất kể mức độ tự động hoá cụ thể nào, kiến trúc phần mềm của một hệ thống lái xe thông minh đều tuân theo một mô hình xử lý ba giai đoạn quen thuộc, thường được gọi là chuỗi cảm nhận -- nhận thức -- lập kế hoạch và điều khiển (sense-perceive-plan-control). Giai đoạn cảm nhận thu thập dữ liệu thô từ các cảm biến gắn trên xe; giai đoạn nhận thức diễn giải dữ liệu thô đó thành các biểu diễn ngữ nghĩa về môi trường xung quanh như vị trí làn đường, biển báo, người đi bộ, phương tiện khác và chướng ngại vật; giai đoạn lập kế hoạch và điều khiển sử dụng kết quả nhận thức để đưa ra quyết định lái và sinh tín hiệu điều khiển truyền tới các cơ cấu chấp hành. Trong ba giai đoạn này, giai đoạn cảm nhận đóng vai trò nền tảng vì mọi quyết định phía sau đều phụ thuộc vào chất lượng và độ tin cậy của dữ liệu cảm biến.

Trong số các cảm biến được sử dụng phổ biến hiện nay như radar, lidar, siêu âm và camera, camera giữ một vị trí đặc biệt quan trọng nhờ khả năng cung cấp thông tin trực quan giàu ngữ nghĩa với chi phí phần cứng thấp. Camera có thể nhận diện màu sắc của đèn tín hiệu, đọc nội dung biển báo, phân biệt người đi bộ với cột đèn hoặc cây xanh, và cung cấp dữ liệu cần thiết cho các mô hình học sâu thực hiện phân vùng ngữ nghĩa (semantic segmentation) cũng như phát hiện đối tượng (object detection). Vì lý do đó, một chiếc xe có khả năng tự lái hiện đại thường được trang bị đồng thời nhiều camera bao quanh thân xe để đảm bảo trường quan sát 360 độ không có điểm mù.

Tuy nhiên, chính việc trang bị nhiều camera lại đặt ra một thách thức rất lớn về băng thông truyền dẫn dữ liệu nội bộ trên xe. Nghiên cứu của Wang và cộng sự cho thấy mỗi camera trong hệ thống AAD có thể sinh ra dòng dữ liệu thô lên tới khoảng 200 MB/s, do đó tám camera đồng thời sẽ tạo ra tổng lưu lượng xấp xỉ 1.600 MB/s \cite{SAC}. Con số này vượt xa khả năng truyền tải của các bus dữ liệu truyền thống trên xe như CAN, FlexRay và thậm chí cả Automotive Ethernet ở các cấu hình phổ biến hiện nay. Việc truyền tải khối lượng dữ liệu khổng lồ này từ các camera tới đơn vị xử lý trung tâm trong thời gian thực, mà không gây nghẽn mạng và không làm tăng quá mức độ trễ hệ thống, đã trở thành một bài toán kỹ thuật mang tính sống còn đối với việc hiện thực hoá phương tiện tự lái ở cấp độ cao.

Trong bối cảnh đó, nén video xuất hiện như một giải pháp tự nhiên để giảm khối lượng dữ liệu truyền tải trên mạng nội bộ của xe. Các chuẩn nén video tiên tiến như H.264/AVC và H.265/HEVC có khả năng giảm kích thước dòng dữ liệu xuống hàng chục đến hàng trăm lần so với dữ liệu thô, đồng thời vẫn giữ được chất lượng hình ảnh đủ tốt cho mắt người quan sát. Tuy nhiên, khác với các kịch bản giải trí truyền thống nơi người dùng cuối là con người, dữ liệu camera trên xe tự lái được tiêu thụ chủ yếu bởi các mô hình trí tuệ nhân tạo. Điều này dẫn đến một nhu cầu nghiên cứu mới về một thế hệ nền tảng nén video không tối ưu cho mắt người mà tối ưu cho thuật toán thị giác máy, đảm bảo rằng dù dòng video bị nén nặng để tiết kiệm băng thông, độ chính xác của các tác vụ nhận thức phía sau vẫn được bảo toàn ở mức tối đa.

\section{Hạn chế của các chuẩn nén video truyền thống trong bài toán thị giác máy}

Các chuẩn nén video phổ biến hiện nay được thiết kế dựa trên đặc tính cảm thụ thị giác của con người. Bộ máy nén thực hiện biến đổi cosin rời rạc (Discrete Cosine Transform), lượng tử hoá các hệ số tần số cao và mã hoá entropy nhằm loại bỏ những thành phần mà mắt người ít nhạy cảm, từ đó đạt được tỉ lệ nén cao mà người xem vẫn cảm thấy hình ảnh chấp nhận được. Toàn bộ chiến lược này hoạt động hiệu quả trong các ứng dụng giải trí, hội nghị truyền hình hay phát trực tuyến, nơi mục tiêu cuối cùng là duy trì trải nghiệm xem cho người dùng.

Tuy nhiên, khi đặt vào bài toán thị giác máy, các giả định nền tảng của những chuẩn nén này lại trở nên không còn phù hợp. Hệ thống mã hoá truyền thống đối xử với mọi vùng ảnh một cách đồng nhất; mỗi macroblock được nén theo cùng một bộ tham số chất lượng mà không quan tâm tới nội dung ngữ nghĩa bên trong nó. Một vùng bầu trời đồng đều và một vùng chứa người đi bộ ở khoảng cách gần đều bị áp cùng một mức độ lượng tử hoá nếu chúng có cùng độ phức tạp tín hiệu. Cách tiếp cận này hợp lý với mắt người vì người xem có khả năng tự bù đắp những chi tiết bị mất, nhưng lại gây thiệt hại đáng kể cho các mô hình học sâu phía sau, vốn dựa vào những chi tiết tinh tế ở biên đối tượng để đưa ra dự đoán chính xác.

Các thí nghiệm được trình bày trong nghiên cứu SAC đã chỉ rõ vấn đề này, khi cho thấy độ chính xác của các mô hình phát hiện đối tượng và phân vùng ngữ nghĩa suy giảm đáng kể khi đầu vào bị nén bằng các chuẩn truyền thống ở tỉ lệ nén cao \cite{SAC}. Các thành phần méo nén như hiệu ứng khối, hiệu ứng vòng, làm mờ biên và mất chi tiết đều có ảnh hưởng trực tiếp tới chất lượng đặc trưng mà mô hình trích xuất được. Đặc biệt, biên của các đối tượng nhỏ ở xa như người đi bộ, biển báo và cọc tiêu chính là những vùng bị tổn thương nhiều nhất sau khi nén, trong khi đây lại là chính những đối tượng an toàn quan trọng mà hệ thống tự lái cần nhận diện sớm và chính xác.

Ngoài ra, các tiêu chí đánh giá chất lượng truyền thống như tỉ số tín hiệu trên nhiễu cực đại (Peak Signal-to-Noise Ratio) và chỉ số tương đồng cấu trúc (Structural Similarity Index) cũng được xây dựng dựa trên cảm thụ của con người, vì vậy chúng không phản ánh đầy đủ chất lượng từ góc nhìn của một mô hình thị giác máy. Một dòng video có thể có chỉ số PSNR cao đối với người xem, nhưng vẫn làm giảm độ chính xác của mô hình phân vùng nếu vùng bị suy giảm chất lượng trùng với vùng ngữ nghĩa quan trọng. Ngược lại, một dòng video có PSNR thấp hơn nhưng giữ nguyên chất lượng ở vùng quan trọng vẫn có thể cho kết quả nhận thức tốt hơn. Điều này cho thấy bản thân các tiêu chí đánh giá cũng cần được suy xét lại khi mục tiêu cuối cùng là phục vụ thuật toán chứ không phải mắt người.

Xuất phát từ những hạn chế này, cộng đồng nghiên cứu đã đặt ra yêu cầu xây dựng một thế hệ nền tảng nén video mới có khả năng phân biệt được mức độ quan trọng ngữ nghĩa giữa các vùng của khung hình, từ đó phân bổ ngân sách bit một cách hợp lý hơn để vừa duy trì băng thông thấp vừa bảo vệ được chất lượng tại những vùng có giá trị nhận thức cao.

\section{Hướng tiếp cận mã hoá video nhận biết ngữ nghĩa}

Một trong những hướng nghiên cứu nổi bật nhằm khắc phục các hạn chế nêu trên là phát triển hệ thống mã hoá video nhận biết ngữ nghĩa, trong đó tiêu biểu là khung làm việc SAC (Semantic-Aware Compression) do Wang và cộng sự đề xuất \cite{SAC}. Ý tưởng cốt lõi của hướng tiếp cận này là sử dụng một mô hình phân vùng ngữ nghĩa chạy trước khi nén để xác định những vùng quan trọng đối với nhiệm vụ thị giác máy, sau đó áp dụng mức độ nén thấp hơn cho các vùng này và mức độ nén cao hơn cho các vùng nền ít quan trọng. Bằng cách phân bổ ngân sách bit không đồng đều theo ngữ nghĩa, hệ thống có thể giữ nguyên chất lượng đặc trưng ở những nơi cần thiết trong khi vẫn đạt được tỉ lệ nén tổng thể tương đương hoặc tốt hơn so với chuẩn truyền thống.

Tuy nhiên, để hướng tiếp cận này khả thi trong môi trường ô tô thời gian thực, mô hình phân vùng ngữ nghĩa phải đáp ứng đồng thời hai yêu cầu khắt khe. Thứ nhất, độ chính xác phải đủ cao để mặt nạ vùng quan tâm phản ánh đúng phân bố đối tượng an toàn trong khung hình, vì một mặt nạ sai sẽ dẫn đến việc nén nặng các vùng quan trọng và ngược lại. Thứ hai, tốc độ suy luận phải đủ nhanh để không trở thành nút thắt cổ chai trong toàn bộ chuỗi xử lý từ camera tới đơn vị nhận thức trung tâm. Một mô hình phân vùng có độ chính xác cao nhưng chạy chậm sẽ làm tăng độ trễ tổng thể và làm mất đi lợi ích về băng thông mà cơ chế nén nhận biết ngữ nghĩa mang lại.

Từ thực tế trên, mạng phân vùng ngữ nghĩa thời gian thực PIDNet do Xu và cộng sự đề xuất xuất hiện như một ứng cử viên rất phù hợp để đóng vai trò xương sống cho nền tảng nén nhận biết ngữ nghĩa \cite{PID}. Bằng cách lấy cảm hứng từ lý thuyết bộ điều khiển PID trong điều khiển tự động, PIDNet sử dụng kiến trúc ba nhánh kết hợp giữa nhánh chi tiết, nhánh ngữ cảnh và nhánh biên để đạt được sự cân bằng tinh tế giữa độ chính xác và tốc độ. Phiên bản PIDNet-S đạt 78,6\% mIOU trên Cityscapes ở tốc độ 93,2 khung hình mỗi giây, còn phiên bản PIDNet-L đạt tới 80,6\% mIOU, dẫn đầu trong nhóm các mô hình phân vùng thời gian thực hiện nay.

Việc tích hợp một mạng phân vùng tiên tiến như PIDNet vào một bộ mã hoá video chuẩn công nghiệp như H.265 đặt ra nhiều thách thức kỹ thuật cần được giải quyết. Mặt nạ vùng quan tâm (Region of Interest) sinh ra ở cấp độ pixel phải được căn chỉnh theo cấu trúc macroblock 16$\times$16 của bộ mã hoá để tránh hiện tượng rò chất lượng giữa các vùng có tham số chất lượng khác nhau. Hai dòng dữ liệu vùng quan tâm và phi vùng quan tâm phải được mã hoá riêng và sau đó kết hợp lại một cách không mất mát. Toàn bộ chuỗi xử lý phải có độ trễ chấp nhận được để có thể triển khai trên xe thật. Đây chính là không gian nghiên cứu và thực nghiệm mà đồ án này hướng tới.

\section{Mục tiêu, phạm vi và đóng góp của đồ án}

Xuất phát từ những phân tích trên, mục tiêu chính của đồ án này là nghiên cứu, thiết kế và hiện thực hoá một nền tảng mã hoá video nhận biết ngữ nghĩa hoàn chỉnh dành cho các ứng dụng thị giác máy, đặc biệt là trong bối cảnh hệ thống camera trên xe tự lái. Nền tảng được xây dựng dựa trên việc kết hợp mạng phân vùng ngữ nghĩa thời gian thực PIDNet với chuẩn nén video tiên tiến H.265/HEVC, đồng thời áp dụng cơ chế phân tách dòng dữ liệu theo vùng quan tâm như đề xuất trong khung SAC.

Về phạm vi, đồ án giới hạn việc thí nghiệm trên tập dữ liệu Cityscapes, vốn là bộ dữ liệu phân vùng ngữ nghĩa lái xe đô thị tiêu chuẩn được sử dụng phổ biến trong cộng đồng nghiên cứu thị giác máy tính. Sơ đồ phân vùng được rút gọn về bốn lớp ngữ nghĩa chính phù hợp với bài toán nén nhận biết: vùng quan tâm bao gồm đường, phương tiện, người và các đối tượng động; ba lớp phi vùng quan tâm gồm bầu trời, công trình xây dựng và thảm thực vật. Việc rút gọn này phản ánh đúng yêu cầu thực tế của bài toán nén, vốn không cần phân biệt chi tiết các lớp con bên trong nhóm vùng quan tâm hay nhóm phi vùng quan tâm. Bộ mã hoá được lựa chọn là x265, một hiện thực mã nguồn mở của chuẩn H.265, do tính phổ biến, hiệu năng cao và khả năng tuỳ biến tham số chất lượng theo từng dòng dữ liệu riêng biệt.

Đồ án đặt ra một số đóng góp cụ thể. Trước hết, đồ án hiện thực hoá thành công một nền tảng phần mềm hoàn chỉnh, kết nối từ khâu huấn luyện mạng phân vùng, sinh mặt nạ vùng quan tâm trên từng khung hình, tách dòng dữ liệu theo vùng, mã hoá song song với hai mức hệ số tốc độ không đổi (Constant Rate Factor) khác nhau, đến khâu kết hợp dòng và đánh giá chất lượng đầu ra. Tiếp theo, đồ án triển khai và áp dụng các tiêu chí đánh giá chất lượng nhận biết ngữ nghĩa SA-PSNR và SA-SSIM được đề xuất trong nghiên cứu SAC \cite{SAC}, cho phép so sánh nền tảng đề xuất với nén truyền thống theo cách phản ánh đúng nhu cầu của các thuật toán thị giác máy. Đồ án cũng thực hiện đánh giá độ trễ từng giai đoạn của chuỗi xử lý trên cả GPU và CPU để khảo sát khả năng triển khai thực tế. Cuối cùng, đồ án trình bày một khảo sát thực nghiệm về mặt cân bằng tốc độ và chất lượng giữa các phiên bản mạng phân vùng khác nhau, từ đó đưa ra các khuyến nghị về cấu hình phù hợp.

Phần còn lại của báo cáo được tổ chức như sau. Chương 2 trình bày cơ sở lý thuyết về nén video, phân vùng ngữ nghĩa, khái niệm vùng quan tâm trong mã hoá và mạng PIDNet. Chương 3 mô tả chi tiết thiết kế hệ thống và quy trình hiện thực hoá nền tảng. Chương 4 trình bày kết quả thực nghiệm về chất lượng nén, độ chính xác phân vùng, các tiêu chí nhận biết ngữ nghĩa và độ trễ. Chương 5 tổng kết các kết quả đạt được và đề xuất các hướng nghiên cứu tiếp theo.


% ============================================================
\chapter{Cơ sở lý thuyết}
% ============================================================

\section{Nén video và các chuẩn hiện đại}

Nén video là quá trình biểu diễn một chuỗi khung hình bằng một dòng bit có kích thước nhỏ hơn nhiều so với dữ liệu thô ban đầu, sao cho khi giải nén ta có thể phục hồi lại chuỗi khung hình tương đương về mặt cảm thụ. Để đạt được điều này, các bộ mã hoá hiện đại khai thác triệt để hai loại dư thừa cơ bản tồn tại trong video tự nhiên. Loại dư thừa thứ nhất là dư thừa không gian, xuất hiện do các pixel lân cận trong cùng một khung hình thường có giá trị tương quan cao với nhau, đặc biệt là trong các vùng bề mặt đồng nhất như bầu trời hay tường nhà. Loại dư thừa thứ hai là dư thừa thời gian, xuất hiện do các khung hình liên tiếp trong một cảnh quay thường rất giống nhau ngoại trừ các vùng có chuyển động.

Đối với dư thừa không gian, các bộ mã hoá áp dụng kỹ thuật mã hoá nội khung (intra-frame coding), trong đó khung hình được chia thành các khối nhỏ và mỗi khối được biến đổi sang miền tần số thông qua biến đổi cosin rời rạc. Sau biến đổi, năng lượng tín hiệu được tập trung chủ yếu ở các hệ số tần số thấp, trong khi các hệ số tần số cao thường có giá trị nhỏ và đại diện cho chi tiết tinh tế. Bước tiếp theo là lượng tử hoá, trong đó các hệ số được chia cho một bước lượng tử rồi làm tròn về số nguyên gần nhất. Đây chính là bước duy nhất mang tính mất mát trong toàn bộ chuỗi xử lý, vì việc làm tròn không thể đảo ngược hoàn toàn. Cuối cùng, dòng hệ số đã lượng tử hoá được nén tiếp bằng mã hoá entropy (entropy coding), ví dụ như mã hoá độ dài thay đổi thích nghi theo ngữ cảnh, để khai thác phân bố thống kê không đều của các giá trị.

Đối với dư thừa thời gian, các bộ mã hoá áp dụng kỹ thuật mã hoá liên khung dựa trên ước lượng chuyển động và bù chuyển động (motion estimation and compensation). Mỗi khối trong khung hình hiện tại được dự đoán từ một khối tham chiếu phù hợp trong một khung hình đã giải mã trước đó hoặc sau đó, vị trí của khối tham chiếu được mô tả bằng một vector chuyển động. Bộ mã hoá chỉ truyền vector chuyển động cùng với phần dư dự đoán (prediction residual), tức là hiệu giữa khối thực và khối dự đoán, thay vì truyền toàn bộ giá trị pixel. Vì phần dư dự đoán thường có biên độ nhỏ hơn nhiều so với pixel gốc, sau khi áp dụng cùng chuỗi biến đổi và lượng tử hoá, lượng bit cần thiết để mã hoá nó giảm đi đáng kể.

Sự đánh đổi giữa kích thước dòng bit và chất lượng hình ảnh được mô tả bởi lý thuyết tốc độ -- méo nén (rate-distortion theory). Lý thuyết này khẳng định rằng với một mức méo cho phép cho trước, tồn tại một giới hạn dưới về tốc độ bit cần thiết để biểu diễn nguồn tín hiệu, và mọi bộ mã hoá thực tế đều cố gắng tiến gần tới giới hạn lý thuyết này. Trong thực tế, các bộ mã hoá hiện đại không tối ưu trực tiếp ràng buộc nén -- chất lượng mà tối ưu một hàm Lagrange dạng
\begin{equation}
  J = D + \lambda R,
\end{equation}
trong đó $D$ là độ méo, $R$ là tốc độ bit và $\lambda$ là hệ số nhân điều khiển sự cân bằng giữa hai mục tiêu.

Hai chuẩn nén video phổ biến nhất hiện nay là H.264/AVC và H.265/HEVC. Cả hai chuẩn đều dựa trên kiến trúc mã hoá lai ghép kết hợp dự đoán nội khung, dự đoán liên khung, biến đổi và lượng tử hoá, nhưng H.265 cải tiến đáng kể về nhiều mặt \cite{SAC}. H.264 sử dụng đơn vị macroblock cố định kích thước 16$\times$16 pixel, trong khi H.265 mở rộng sang đơn vị mã hoá phân cấp với kích thước thay đổi từ 8$\times$8 đến 64$\times$64 pixel, giúp tăng hiệu quả nén trên các vùng phẳng lớn cũng như các vùng chi tiết nhỏ. H.265 cũng hỗ trợ nhiều chế độ dự đoán nội khung hơn, ước lượng chuyển động chính xác hơn và bộ lọc khử khối được cải tiến. Tổng hợp lại, H.265 có thể giảm khoảng 50\% tốc độ bit so với H.264 ở cùng mức chất lượng cảm thụ.

Một tham số điều khiển chất lượng quan trọng được sử dụng rộng rãi trong các bộ mã hoá hiện thực hai chuẩn trên là CRF. Tham số CRF nằm trong khoảng từ 0 đến 51 và điều khiển mức độ lượng tử hoá một cách đồng đều trên toàn bộ khung hình. Giá trị CRF càng nhỏ thì độ lượng tử hoá càng thấp, dẫn đến chất lượng cao hơn và dòng bit lớn hơn; ngược lại, CRF càng lớn thì độ nén càng cao nhưng chất lượng càng giảm. Một mối quan hệ xấp xỉ thường gặp giữa CRF và bước lượng tử thực tế có thể được mô tả qua công thức
\begin{equation}
  Q_{\text{step}} = 2^{(\text{CRF} - 4)/6}.
\end{equation}
Điểm cần lưu ý là CRF chỉ áp dụng một mức chất lượng duy nhất cho cả khung hình, không phân biệt vùng quan trọng và vùng nền; đây chính là hạn chế nền tảng mà các phương pháp nén nhận biết ngữ nghĩa hướng tới khắc phục.

Cuối cùng, để đánh giá chất lượng dòng video sau nén, hai tiêu chí được sử dụng phổ biến là PSNR và SSIM. PSNR đo lường sai số bình phương trung bình giữa khung gốc và khung tái tạo trên miền pixel, được tính qua giá trị pixel cực đại $L$ và sai số bình phương trung bình $\text{MSE}$:
\begin{equation}
  \text{PSNR} = 10 \log_{10} \frac{L^2}{\text{MSE}}.
\end{equation}
SSIM thì tập trung vào sự tương đồng cấu trúc giữa hai ảnh thông qua so sánh cường độ trung bình, độ tương phản và cấu trúc cục bộ. Cả hai tiêu chí này đều được thiết kế dựa trên cảm thụ thị giác của con người và đối xử với mọi vùng pixel trong khung hình một cách bình đẳng, do đó không phản ánh đầy đủ chất lượng từ góc nhìn của một mô hình thị giác máy, như đã phân tích ở chương trước.

\section{Phân vùng ngữ nghĩa trong thị giác máy}

Phân vùng ngữ nghĩa là một trong những bài toán cơ bản và quan trọng nhất của thị giác máy tính (computer vision), trong đó mục tiêu là gán cho mỗi pixel của ảnh đầu vào một nhãn lớp ngữ nghĩa thuộc một tập lớp cho trước. Khác với bài toán phân loại ảnh chỉ trả về một nhãn duy nhất cho toàn ảnh hay bài toán phát hiện đối tượng chỉ trả về các hộp bao quanh, phân vùng ngữ nghĩa cung cấp một biểu diễn dày đặc ở cấp độ pixel, cho phép mô tả chi tiết hình dạng và ranh giới của từng đối tượng trong cảnh. Tính chất dày đặc này biến phân vùng ngữ nghĩa trở thành một công cụ không thể thiếu trong các ứng dụng đòi hỏi hiểu cảnh chính xác như xe tự lái, robot tự hành và phân tích ảnh y tế.

Trong bối cảnh ô tô, phân vùng ngữ nghĩa giữ vai trò đặc biệt quan trọng vì nó cung cấp bản đồ ngữ nghĩa phong phú mà các tầng nhận thức cao hơn có thể khai thác. Một mô hình phân vùng tốt cho phép xác định chính xác vị trí của làn đường, vạch kẻ, vỉa hè, phương tiện khác, người đi bộ, biển báo, cọc tiêu, công trình hai bên đường, thảm thực vật và bầu trời. Trên cơ sở bản đồ này, các thuật toán phía sau có thể quyết định hành vi lái hợp lý, xác định vùng có thể di chuyển an toàn và tránh va chạm với các đối tượng động. Vì lý do này, ngành công nghiệp ô tô đã đầu tư mạnh vào việc xây dựng các tập dữ liệu phân vùng quy mô lớn như Cityscapes, BDD100K hay KITTI-STEP, đồng thời thúc đẩy nhiều cải tiến kiến trúc mạng học sâu chuyên biệt.

Sự phát triển của các kiến trúc phân vùng ngữ nghĩa đã trải qua nhiều thế hệ. Khởi đầu là mạng tích chập đầy đủ (Fully Convolutional Network), trong đó các lớp kết nối đầy đủ ở cuối mạng phân loại được thay bằng các lớp tích chập để cho phép đầu ra có cấu trúc không gian. Tiếp đó là các kiến trúc mã hoá -- giải mã (encoder-decoder) như U-Net và SegNet, sử dụng đường nối tắt từ tầng mã hoá sang tầng giải mã để khôi phục chi tiết không gian bị mất do giảm mẫu. Một bước tiến quan trọng kế tiếp là các phương pháp sử dụng cơ chế chú ý như nhóm DeepLab với tích chập giãn (atrous convolution) và pooling theo nhiều tỉ lệ, cho phép mở rộng trường tiếp nhận mà không làm giảm độ phân giải đặc trưng. Gần đây hơn, các kiến trúc hai nhánh như BiSeNet, DDRNet và STDC chia mạng thành một nhánh ngữ cảnh sâu để học biểu diễn toàn cục và một nhánh chi tiết nông để bảo toàn thông tin biên, sau đó hợp nhất hai nhánh ở cuối để cho ra dự đoán cuối cùng.

Tiêu chí đánh giá phổ biến nhất cho phân vùng ngữ nghĩa là chỉ số trung bình giao trên hợp (mean Intersection over Union). Với mỗi lớp ngữ nghĩa, ta tính tỉ số giữa diện tích giao và diện tích hợp của vùng dự đoán và vùng tham chiếu, sau đó lấy trung bình trên toàn bộ các lớp. Công thức tổng quát của mIOU có thể được viết như sau:
\begin{equation}
  \text{mIOU} = \frac{1}{m} \sum_{c=1}^{m} \frac{|\text{GT}_c \cap \text{Pre}_c|}{|\text{GT}_c \cup \text{Pre}_c|},
\end{equation}
trong đó $m$ là số lớp, $\text{GT}_c$ và $\text{Pre}_c$ lần lượt là tập pixel thuộc lớp $c$ theo nhãn chuẩn và theo dự đoán của mô hình. Chỉ số mIOU phản ánh đồng thời sự chính xác của dự đoán dương và sự bao phủ của dự đoán đối với toàn bộ đối tượng, đồng thời không thiên vị về phía các lớp có diện tích lớn vì được lấy trung bình theo lớp.

Tuy nhiên, mặc dù mIOU là một chỉ số hữu ích, các nhà nghiên cứu phải đối mặt với nhiều thách thức khi triển khai phân vùng ngữ nghĩa cho ứng dụng thực tế. Thách thức đầu tiên là sự đánh đổi giữa độ chính xác và tốc độ; các mô hình có mIOU cao thường có kích thước lớn và chi phí tính toán cao, không phù hợp với các thiết bị có tài nguyên hạn chế trên xe. Thách thức thứ hai là độ chính xác tại biên đối tượng, đây là vùng có ý nghĩa quyết định đối với các tác vụ phía sau nhưng cũng là vùng mà nhiều mô hình mắc lỗi nhiều nhất do hiện tượng làm mờ biên qua các tầng giảm mẫu. Thách thức thứ ba là phát hiện đối tượng nhỏ ở khoảng cách xa, vốn chiếm một tỉ lệ pixel rất nhỏ trong khung hình nhưng có thể là người đi bộ hay xe đạp cần tránh va chạm. Những thách thức này đã thúc đẩy việc phát triển các kiến trúc mạng tiên tiến hơn, trong đó hướng tiếp cận hai nhánh và đa nhánh được xem là một con đường đầy hứa hẹn.

Cụ thể, kiến trúc mạng hai nhánh chia luồng đặc trưng thành hai phần song song. Nhánh ngữ cảnh đi sâu vào mạng, thực hiện nhiều bước giảm mẫu để có trường tiếp nhận lớn và học được mối quan hệ phụ thuộc tầm xa giữa các vùng trong ảnh. Nhánh chi tiết giữ độ phân giải cao, đi qua ít tầng giảm mẫu hơn, tập trung vào việc bảo toàn thông tin biên và cấu trúc cục bộ. Hai nhánh sau đó được hợp nhất thông qua một mô-đun hợp nhất đặc trưng (feature fusion module) có thể huấn luyện được. Phương pháp này đã đạt được sự cân bằng tốt giữa độ chính xác và tốc độ, tuy nhiên nó vẫn để lại một vấn đề mở liên quan đến hiện tượng vượt mức tại biên đối tượng, vấn đề này sẽ được phân tích chi tiết hơn ở phần sau khi trình bày về mạng PIDNet.

\section{Vùng quan tâm trong mã hoá video}

Khái niệm vùng quan tâm đã xuất hiện từ lâu trong lĩnh vực xử lý ảnh và mã hoá video. Trong ngữ cảnh nén ảnh JPEG2000, vùng quan tâm cho phép người dùng chỉ định một vùng được nén với chất lượng cao hơn so với phần còn lại của ảnh. Trong các hệ thống hội nghị truyền hình truyền thống, vùng chứa khuôn mặt người nói thường được xác định và được phân bổ nhiều bit hơn để duy trì biểu cảm rõ ràng, trong khi vùng nền cố định có thể bị nén mạnh mà không gây ảnh hưởng đáng kể tới trải nghiệm. Một số hiện thực phần mềm mã hoá video thậm chí cung cấp tính năng ánh xạ chất lượng theo vùng, cho phép người dùng cấp một bản đồ chất lượng tuỳ chỉnh để điều khiển bộ mã hoá.

Tuy nhiên, các cách tiếp cận vùng quan tâm truyền thống đều mang tính tĩnh hoặc bán tĩnh, nghĩa là vùng quan tâm được xác định trước hoặc thay đổi rất chậm theo thời gian, và quan trọng hơn, chúng được thiết kế cho mắt người chứ không cho thị giác máy. Khi áp dụng vào bài toán camera trên xe tự lái, cách tiếp cận này không còn phù hợp vì vùng quan tâm thực sự không cố định mà thay đổi liên tục theo nội dung động của cảnh: ở mỗi khung hình, vị trí của xe, người và đối tượng động đều khác nhau, đồng thời mức độ quan trọng của từng vùng cũng phụ thuộc vào ngữ cảnh lái cụ thể.

Trong bối cảnh đó, khung làm việc SAC đề xuất sử dụng một mô hình học máy chạy trước khâu nén để sinh ra mặt nạ vùng quan tâm động cho từng khung hình \cite{SAC}. Mặt nạ này phân loại các pixel của khung hình thành hai nhóm: vùng quan tâm bao gồm đường, phương tiện, người đi bộ và các đối tượng an toàn quan trọng khác; phi vùng quan tâm bao gồm các thành phần nền như bầu trời, công trình xây dựng và thảm thực vật. Cách phân chia này phản ánh đúng mức độ quan trọng của các thành phần khác nhau trong khung hình đối với hệ thống lái tự động: việc nhận diện sai một người đi bộ có thể dẫn đến tai nạn nghiêm trọng, trong khi việc làm giảm chi tiết của một mảng bầu trời hầu như không ảnh hưởng đến quyết định lái.

Một vấn đề kỹ thuật quan trọng phát sinh khi áp dụng mặt nạ ở cấp độ pixel vào bộ mã hoá video làm việc trên đơn vị macroblock. Vì các bộ mã hoá H.264 và H.265 thực hiện biến đổi và lượng tử hoá trên các khối cố định 16$\times$16 pixel, nếu biên của mặt nạ vùng quan tâm không trùng với biên macroblock, sẽ xuất hiện những khối nằm đè lên cả vùng quan tâm và phi vùng quan tâm. Khi mã hoá hai dòng dữ liệu riêng biệt, các khối hỗn hợp này sẽ chứa cả thông tin từ hai vùng và gây hiện tượng rò chất lượng, làm phá vỡ tính tách biệt giữa hai dòng dữ liệu. Để khắc phục, SAC sử dụng một bộ lọc căn chỉnh macroblock, trong đó một macroblock được gán hoàn toàn cho vùng quan tâm nếu nó chứa bất kỳ pixel quan tâm nào, ngược lại được gán cho phi vùng quan tâm. Cách căn chỉnh này đảm bảo biên mặt nạ luôn trùng với biên macroblock và loại bỏ hoàn toàn hiện tượng rò chất lượng giữa hai dòng \cite{SAC}.

Sau khi mặt nạ đã được căn chỉnh, mỗi khung hình được tách thành hai khung con: khung con vùng quan tâm chỉ chứa giá trị pixel tại các vị trí vùng quan tâm và bằng không ở các vị trí khác, khung con phi vùng quan tâm thì ngược lại. Hai chuỗi khung con này được mã hoá độc lập với hai tham số CRF khác nhau: chuỗi vùng quan tâm được nén ở CRF thấp để giữ chất lượng cao, chuỗi phi vùng quan tâm được nén ở CRF cao để tiết kiệm dòng bit. Sau khi giải mã, hai khung con được kết hợp lại thành khung hình hoàn chỉnh thông qua phép cộng pixel, vì vùng giá trị khác không của hai khung con không chồng lấn nhau. Một điểm cần lưu ý là trong quá trình mã hoá, cơ chế lượng tử thích nghi được tắt hoàn toàn vì việc kiểm soát chất lượng theo vùng đã được thực hiện ở mức tách dòng, và việc bật thêm lượng tử thích nghi sẽ phá vỡ sự thống nhất chất lượng trong nội bộ mỗi dòng.

Để đánh giá chất lượng đầu ra của hệ thống nén nhận biết ngữ nghĩa một cách công bằng và phản ánh đúng nhu cầu của thị giác máy, nghiên cứu SAC đề xuất hai tiêu chí mới là SA-PSNR và SA-SSIM, được định nghĩa như sau \cite{SAC}:
\begin{equation}
  \text{SA-PSNR} = r_{\text{non}} \cdot P_i + r_{\text{roi}} \cdot P_n,
\end{equation}
\begin{equation}
  \text{SA-SSIM} = r_{\text{non}} \cdot S_i + r_{\text{roi}} \cdot S_n,
\end{equation}
trong đó $P_i$ và $S_i$ là PSNR và SSIM tính trên vùng quan tâm, $P_n$ và $S_n$ là PSNR và SSIM tính trên phi vùng quan tâm, còn $r_{\text{roi}}$ và $r_{\text{non}}$ là các chỉ số tỉ lệ nén được xác định bởi
\begin{equation}
  r_{\text{roi}} = \frac{\text{CRF}_{\text{roi}}}{\text{CRF}_{\text{roi}} + \text{CRF}_{\text{non}}}, \qquad
  r_{\text{non}} = \frac{\text{CRF}_{\text{non}}}{\text{CRF}_{\text{roi}} + \text{CRF}_{\text{non}}}.
\end{equation}
Cấu trúc trọng số chéo trong các công thức trên có một ý nghĩa rất tinh tế. Trọng số gắn với chất lượng vùng quan tâm là $r_{\text{non}}$, tức là khi CRF của phi vùng quan tâm càng lớn so với CRF của vùng quan tâm thì chất lượng vùng quan tâm càng được nhân với một hệ số lớn, qua đó khuếch đại đóng góp của vùng quan trọng trong điểm số tổng thể. Cách trọng số hoá này phản ánh đúng triết lý của nén nhận biết ngữ nghĩa: chỉ số đánh giá càng nghiêng về vùng được nén ít hơn, tức là vùng được hệ thống coi trọng hơn về mặt ngữ nghĩa.

Thực nghiệm của Wang và cộng sự cho thấy hệ thống SA-X264 đạt cải thiện trung bình $+2{,}864$~dB SA-PSNR và $+2{,}7\%$ mIOU so với chuẩn H.264 truyền thống ở cùng tỉ lệ nén \cite{SAC}, khẳng định tính đúng đắn của hướng tiếp cận nén theo vùng quan tâm động dựa trên phân vùng ngữ nghĩa.

\section{Mạng PIDNet và lý thuyết bộ điều khiển PID}

Bộ điều khiển PID (Proportional-Integral-Derivative) là một trong những kết quả kinh điển và được sử dụng rộng rãi nhất của lý thuyết điều khiển tự động. Trong sơ đồ điều khiển kín, bộ điều khiển PID tạo ra tín hiệu điều khiển $u(t)$ từ sai số $e(t)$ giữa giá trị đặt và giá trị thực đo, theo công thức
\begin{equation}
  u(t) = K_p \, e(t) + K_i \int_{0}^{t} e(\tau) \, d\tau + K_d \, \frac{de(t)}{dt},
\end{equation}
trong đó ba thành phần thể hiện ba cách phản ứng khác nhau với sai số. Thành phần tỉ lệ $K_p \, e(t)$ phản ứng tức thời với sai số hiện tại; thành phần tích phân $K_i \int e \, d\tau$ tích luỹ sai số trong quá khứ và giúp loại bỏ sai số xác lập; thành phần đạo hàm $K_d \, de/dt$ dự đoán xu hướng thay đổi của sai số trong tương lai dựa trên tốc độ biến thiên, đóng vai trò làm giảm hiện tượng vượt mức (overshoot) và dao động.

Một quan sát quan trọng được Xu và cộng sự đưa ra trong nghiên cứu PIDNet là kiến trúc mạng phân vùng ngữ nghĩa hai nhánh có thể được nhìn nhận như một bộ điều khiển PI trong miền Fourier \cite{PID}. Nhánh chi tiết của mạng hai nhánh thực chất hoạt động như thành phần tỉ lệ, phản ứng nhanh với thông tin tần số cao tại các pixel cục bộ; nhánh ngữ cảnh hoạt động như thành phần tích phân, tích luỹ thông tin tần số thấp trên một vùng rộng để cung cấp ngữ cảnh toàn cục. Tuy nhiên, một bộ điều khiển PI thiếu thành phần đạo hàm, do đó dễ bị hiện tượng vượt mức khi tín hiệu thay đổi nhanh.

Hiện tượng vượt mức trong điều khiển có một biểu hiện rất rõ ràng và gây hại trong bài toán phân vùng ngữ nghĩa. Khi mạng phân vùng đi qua các tầng giảm mẫu để học ngữ cảnh, thông tin tần số thấp được khuếch đại trong khi thông tin tần số cao bị làm mịn. Khi hai loại thông tin này được hợp nhất ở cuối mạng, thành phần tần số thấp lấn át thành phần tần số cao tại các vùng biên đối tượng, khiến mạng dự đoán nhãn ngữ cảnh xâm lấn vào trong biên thực sự của đối tượng. Đây chính là phiên bản không gian của hiện tượng vượt mức trong điều khiển, và nó đặc biệt gây hại cho các đối tượng nhỏ hoặc các đối tượng có hình dạng phức tạp, nơi sai lệch tại biên có thể chiếm tỉ lệ lớn so với toàn bộ đối tượng.

Xuất phát từ những hạn chế này, PIDNet bổ sung thêm một nhánh thứ ba lấy cảm hứng từ thành phần đạo hàm của bộ điều khiển PID, từ đó tạo nên kiến trúc ba nhánh hoàn chỉnh gồm nhánh tỉ lệ (P), nhánh tích phân (I) và nhánh đạo hàm (D) \cite{PID}. Nhánh tỉ lệ là một nhánh nông giữ độ phân giải cao, có nhiệm vụ trích xuất các đặc trưng chi tiết cục bộ giàu thông tin tần số cao, tương ứng với việc phản ứng với sai số hiện tại trong điều khiển. Nhánh tích phân là một nhánh sâu với nhiều tầng giảm mẫu, có nhiệm vụ tổng hợp ngữ cảnh tầm xa thông qua việc mở rộng trường tiếp nhận, đóng vai trò như một bộ lọc thông thấp tương tự thành phần tích phân tích luỹ thông tin trong quá khứ. Nhánh đạo hàm là một nhánh phụ trợ chuyên trích xuất các đặc trưng biên đối tượng, đóng vai trò như một bộ lọc thông cao để bù trừ cho xu hướng làm mịn của nhánh tích phân và giảm thiểu hiện tượng vượt mức tại biên.

Mối liên hệ toán học giữa các nhánh mạng và các thành phần PID có thể được phân tích trong miền tần số. Nếu thay $z^{-1}$ bằng $e^{-j\omega}$ trong biểu diễn miền $z$, hàm truyền của bộ điều khiển PID rời rạc có thể viết thành
\begin{equation}
  C(z) = K_p + K_i\left(1 - e^{-j\omega}\right)^{-1} + K_d\left(1 - e^{-j\omega}\right),
\end{equation}
cho thấy rằng khi tần số $\omega$ tăng, độ khuếch đại của thành phần I giảm đi trong khi độ khuếch đại của thành phần D tăng lên \cite{PID}. Do đó ba thành phần P, I, D lần lượt hoạt động như bộ lọc thông toàn dải, thông thấp và thông cao trong miền tần số, bổ sung cho nhau một cách tự nhiên để xử lý đầy đủ cả thông tin tần số cao và tần số thấp.

Để khai thác hiệu quả thông tin từ ba nhánh, PIDNet sử dụng ba mô-đun hợp nhất được thiết kế chuyên biệt. Mô-đun thứ nhất là Pag (Pixel-attention-guided fusion module), thực hiện hợp nhất đặc trưng từ nhánh P và nhánh I ở giai đoạn giữa của mạng. Pag sử dụng một cổng sigmoid tính từ đặc trưng nhánh I để quyết định mức độ tin cậy của đặc trưng chi tiết từ nhánh P tại từng pixel: nếu ngữ cảnh chỉ ra rằng pixel đang thuộc về một vùng đồng nhất thì cổng sẽ ưu tiên đặc trưng ngữ cảnh, ngược lại nếu pixel nằm gần biên thì cổng sẽ ưu tiên đặc trưng chi tiết. Mô-đun thứ hai là PAPPM (Parallel Aggregation Pyramid Pooling Module), được sử dụng ở cuối nhánh I để tổng hợp đa tỉ lệ ngữ cảnh. Khác với mô-đun PPM (Pyramid Pooling Module) truyền thống có cấu trúc tuần tự, PAPPM tổ chức các nhánh pooling theo dạng song song nhằm giảm độ sâu tính toán và tăng tốc độ suy luận mà vẫn giữ được khả năng tổng hợp thông tin đa tỉ lệ. Mô-đun thứ ba là Bag (Boundary-attention-guided fusion module), thực hiện hợp nhất cuối cùng của ba luồng đặc trưng với sự dẫn hướng của bản đồ biên do nhánh D dự đoán. Cụ thể, tại các pixel mà nhánh D dự đoán xác suất biên cao, mô-đun Bag sẽ ưu tiên đặc trưng chi tiết từ nhánh P; tại các pixel có xác suất biên thấp, mô-đun sẽ ưu tiên đặc trưng ngữ cảnh từ nhánh I. Cơ chế dẫn hướng theo biên này giúp giảm thiểu hiện tượng vượt mức một cách trực tiếp và hiệu quả \cite{PID}.

Quá trình huấn luyện PIDNet sử dụng một hàm mất mát tổng hợp gồm bốn thành phần:
\begin{equation}
  \mathcal{L} = \lambda_0 l_0 + \lambda_1 l_1 + \lambda_2 l_2 + \lambda_3 l_3,
\end{equation}
trong đó $l_0$ là mất mát phân vùng chính tính trên đầu ra cuối cùng của mạng, $l_1$ là mất mát biên áp dụng cho đầu ra của nhánh D, $l_2$ là mất mát phụ trợ tính trên một đầu ra trung gian sau Pag, và $l_3$ là mất mát nhận biết biên (Boundary-Awareness Loss) áp dụng có chọn lọc trên các pixel gần biên. Sự kết hợp của bốn thành phần này tạo nên một tín hiệu giám sát đa khía cạnh, vừa thúc đẩy độ chính xác phân vùng toàn cục, vừa cải thiện chất lượng biên và vừa hỗ trợ học nhanh hơn nhờ giám sát sâu. Các trọng số được đặt theo thực nghiệm là $\lambda_0 = 0{,}4$, $\lambda_1 = 20$, $\lambda_2 = 1$, $\lambda_3 = 1$ \cite{PID}.

Về hiệu năng, PIDNet được phát hành ở ba phiên bản với kích thước tăng dần là PIDNet-S, PIDNet-M và PIDNet-L. Trên tập kiểm thử Cityscapes ở độ phân giải $2048 \times 1024$, PIDNet-S đạt mIOU $78{,}6\%$ ở tốc độ $93{,}2$ khung hình mỗi giây, trong khi PIDNet-L đạt mIOU $80{,}6\%$, là kết quả tốt nhất trong nhóm mô hình phân vùng thời gian thực tại thời điểm công bố \cite{PID}. Sự cân bằng giữa độ chính xác cao và tốc độ suy luận nhanh khiến PIDNet trở thành một lựa chọn rất phù hợp cho vai trò xương sống phân vùng ngữ nghĩa trong nền tảng nén video được đề xuất trong đồ án này. Việc tích hợp PIDNet vào chuỗi xử lý SAC sẽ cho phép sinh mặt nạ vùng quan tâm chính xác trên từng khung hình mà không trở thành nút thắt cổ chai về độ trễ, qua đó hiện thực hoá đầy đủ tiềm năng của hướng tiếp cận mã hoá video nhận biết ngữ nghĩa.
